\documentclass{report}
\usepackage{luatexja} % LuaTeXで日本語を使うためのパッケージ
\usepackage{luatexja-fontspec} % LuaTeX用の日本語フォント設定

% --- 数学関連 ---
\usepackage{amsmath, amssymb, amsfonts, mathtools, bm, amsthm} % 基本的な数学パッケージ
\usepackage{type1cm, upgreek} % 数式フォントとギリシャ文字
\usepackage{physics, mhchem} % 物理や化学の記号や式の表記を簡単にする

% --- 表関連 ---
\usepackage{multirow, longtable, tabularx, array, colortbl, dcolumn, diagbox} % 表のレイアウトを柔軟にする
\usepackage{tablefootnote, truthtable} % 表中に注釈を追加、真理値表
\usepackage{tabularray} % 高度な表組みレイアウト

% --- グラフィック関連 ---
\usepackage{tikz, graphicx} % 図の描画と画像の挿入
\usepackage{background} % ウォーターマークの設定
\usepackage{caption, subcaption} % 図や表のキャプション設定
\usepackage{float, here} % 図や表の位置指定

% --- レイアウトとページ設定 ---
\usepackage{fancyhdr} % ページヘッダー、フッター、余白の設定
\usepackage[top = 20truemm, bottom = 20truemm, left = 20truemm, right = 20truemm]{geometry}
\usepackage{fancybox, ascmac} % ボックスのデザイン

% --- 色とスタイル ---
\usepackage{xcolor, color, colortbl, tcolorbox} % 色とカラーボックス
\usepackage{listings, jvlisting} % コードの色付けとフォーマット

% --- 参考文献関連 ---
\usepackage{biblatex, usebib} % 参考文献の管理と挿入
\usepackage{url, hyperref} % URLとリンクの設定

% --- その他の便利なパッケージ ---
\usepackage{footmisc} % 脚注のカスタマイズ
\usepackage{multicol} % 複数段組
\usepackage{comment} % コメントアウトの拡張
\usepackage{siunitx} % 単位の表記
\usepackage{docmute}
% \usepackage{appendix}
% --- tcolorboxとtikzの設定 ---
\tcbuselibrary{theorems, breakable} % 定理のボックスと改ページ設定
\usetikzlibrary{decorations.markings, arrows.meta, calc} % tikzの装飾や矢印の設定

% --- 定理スタイルと数式設定 ---
\theoremstyle{definition} % 定義スタイル
\numberwithin{equation}{section} % 式番号をサブセクション単位でリセット

% --- hyperrefの設定 ---
\hypersetup{
  setpagesize = false,
  bookmarks = true,
  bookmarksdepth = tocdepth,
  bookmarksnumbered = true,
  colorlinks = false,
  pdftitle = {}, % PDFタイトル
  pdfsubject = {}, % PDFサブジェクト
  pdfauthor = {}, % PDF作者
  pdfkeywords = {} % PDFキーワード
}

% --- siunitxの設定 ---
\sisetup{
  table-format = 1.5, % 小数点以下の桁数
  table-number-alignment = center, % 数値の中央揃え
}

% --- 透かし画像の設定 --- 
\backgroundsetup{
  scale=0.5,                       % 画像のスケール
  % color=black,                   % 画像の色（透かし用に半透明が推奨）
  opacity=0.2,                   % 透かしの透明度（0が完全透明、1が完全不透明）
  angle=0,                       % 画像の角度
  position = current page.south east,  % ページの右下
  hshift=-6cm, % 右方向へのシフト（負の値で内側に移動）
  vshift=5cm,  % 上方向へのシフト（正の値で内側に移動）
  contents={\includegraphics{./fig/appilogo-circular-full.png}} % 画像のパス
}

% --- その他の設定 ---
\allowdisplaybreaks % 数式の途中改ページ許可
\newcolumntype{t}{!{\vrule width 0.1pt}} % 新しいカラムタイプ
\newcolumntype{b}{!{\vrule width 1.5pt}} % 太いカラム
\UseTblrLibrary{amsmath, booktabs, counter, diagbox, functional, hook, html, nameref, siunitx, varwidth, zref} % tabularrayのライブラリ
\setlength{\columnseprule}{0.4pt} % カラム区切り線の太さ
\captionsetup[figure]{font = bf} % 図のキャプションの太字設定
\captionsetup[table]{font = bf} % 表のキャプションの太字設定
\captionsetup[lstlisting]{font = bf} % コードのキャプションの太字設定
\captionsetup[subfigure]{font = bf, labelformat = simple} % サブ図のキャプション設定
\setcounter{secnumdepth}{4} % セクションの深さ設定
\newcolumntype{d}{D{.}{.}{5}} % 数値のカラム
\newcolumntype{M}[1]{>{\centering\arraybackslash}m{#1}} % センター揃えのカラム
\DeclareMathOperator{\diag}{diag}
\everymath{\displaystyle} % 数式のスタイル
\newcommand{\inner}[2]{\left\langle #1, #2 \right\rangle}
\renewcommand{\figurename}{図}
\renewcommand{\i}{\mathrm{i}} % 複素数単位i
\renewcommand{\laplacian}{\grad^2} % ラプラシアンの記号
\renewcommand{\thesubfigure}{(\alph{subfigure})} % サブ図の番号形式
\newcommand{\m}[3]{\multicolumn{#1}{#2}{#3}} % マルチカラムのショートカット
\renewcommand{\r}[1]{\mathrm{#1}} % mathrmのショートカット
\newcommand{\e}{\mathrm{e}} % 自然対数の底e
\newcommand{\Ef}{E_{\mathrm{F}}} % フェルミエネルギー
\renewcommand{\c}{\si{\degreeCelsius}} % 摂氏記号
\renewcommand{\d}{\r{d}} % d記号
\renewcommand{\t}[1]{\texttt{#1}} % タイプライタフォント
\newcommand{\kb}{k_{\mathrm{B}}} % ボルツマン定数
% \renewcommand{\phi}{\varphi} % ϕをφに変更
\renewcommand{\epsilon}{\varepsilon}
\newcommand{\fullref}[1]{\textbf{\ref{#1} \nameref{#1}}}
\newcommand{\reff}[1]{\textbf{図\ref{#1}}} % 図参照のショートカット
\newcommand{\reft}[1]{\textbf{表\ref{#1}}} % 表参照のショートカット
\newcommand{\refe}[1]{\textbf{式\eqref{#1}}} % 式参照のショートカット
\newcommand{\refp}[1]{\textbf{コード\ref{#1}}} % コード参照のショートカット
\renewcommand{\lstlistingname}{コード} % コードリストの名前
\renewcommand{\theequation}{\thesection.\arabic{equation}} % 式番号の形式
\renewcommand{\footrulewidth}{0.4pt} % フッターの線
\newcommand{\mar}[1]{\textcircled{\scriptsize #1}} % 丸囲み文字
\newcommand{\combination}[2]{{}_{#1} \mathrm{C}_{#2}} % 組み合わせ
\newcommand{\thline}{\noalign{\hrule height 0.1pt}} % 細い横線
\newcommand{\bhline}{\noalign{\hrule height 1.5pt}} % 太い横線

% --- カスタム色定義 ---
\definecolor{burgundy}{rgb}{0.5, 0.0, 0.13} % バーガンディ色
\definecolor{charcoal}{rgb}{0.21, 0.27, 0.31} % チャコール色
\definecolor{forest}{rgb}{0.0, 0.35, 0} % 森の緑色

% --- カスタム定理環境の定義 ---
\newtcbtheorem[number within = chapter]{myexc}{練習問題}{
  fonttitle = \gtfamily\sffamily\bfseries\upshape,
  colframe = forest,
  colback = forest!2!white,
  rightrule = 1pt,
  leftrule = 1pt,
  bottomrule = 2pt,
  colbacktitle = forest,
  theorem style = standard,
  breakable,
  arc = 0pt,
}{exc-ref}
\newtcbtheorem[number within = chapter]{myprop}{命題}{
  fonttitle = \gtfamily\sffamily\bfseries\upshape,
  colframe = blue!50!black,
  colback = blue!50!black!2!white,
  rightrule = 1pt,
  leftrule = 1pt,
  bottomrule = 2pt,
  colbacktitle = blue!50!black,
  theorem style = standard,
  breakable,
  arc = 0pt
}{proposition-ref}
\newtcbtheorem[number within = chapter]{myrem}{注意}{
  fonttitle = \gtfamily\sffamily\bfseries\upshape,
  colframe = yellow!20!black,
  colback = yellow!50,
  rightrule = 1pt,
  leftrule = 1pt,
  bottomrule = 2pt,
  colbacktitle = yellow!20!black,
  theorem style = standard,
  breakable,
  arc = 0pt
}{remark-ref}
\newtcbtheorem[number within = chapter]{myex}{例題}{
  fonttitle = \gtfamily\sffamily\bfseries\upshape,
  colframe = black,
  colback = white,
  rightrule = 1pt,
  leftrule = 1pt,
  bottomrule = 2pt,
  colbacktitle = black,
  theorem style = standard,
  breakable,
  arc = 0pt
}{example-ref}
\newtcbtheorem[number within = chapter]{exc}{Requirement}{myexc}{exc-ref}
\newcommand{\rqref}[1]{{\bfseries\sffamily 練習問題 \ref{exc-ref:#1}}}
\newtcbtheorem[number within = chapter]{definition}{Definition}{mydef}{definition-ref}
\newcommand{\dfref}[1]{{\bfseries\sffamily 定義 \ref{definition-ref:#1}}}
\newtcbtheorem[number within = chapter]{prop}{命題}{myprop}{proposition-ref}
\newcommand{\prref}[1]{{\bfseries\sffamily 命題 \ref{proposition-ref:#1}}}
\newtcbtheorem[number within = chapter]{rem}{注意}{myrem}{remark-ref}
\newcommand{\rmref}[1]{{\bfseries\sffamily 注意 \ref{remark-ref:#1}}}
\newtcbtheorem[number within = chapter]{ex}{例題}{myex}{example-ref}
\newcommand{\exref}[1]{{\bfseries\sffamily 例題 \ref{example-ref:#1}}}
% --- 再定義コマンド ---
% \mathtoolsset{showonlyrefs=true} % 必要な式番号のみ表示
\pagestyle{fancy} % ヘッダー・フッターのスタイル設定
\chead{応用量子物性講義ノート} % 中央ヘッダー
% \rhead{}
\fancyhead[R]{\rightmark}
\renewcommand{\sectionmark}[1]{\markright{\thesection\ #1}}
\cfoot{\thepage} % 中央フッターにページ番号
\lhead{}
\rfoot{Yuto Masuda and Haruki Aoki} % 右フッターに名前
\setcounter{tocdepth}{4} % 目次の深さ
\makeatletter
\@addtoreset{equation}{section} % サブセクションごとに式番号をリセット
\makeatother

% --- メタ情報 ---
\title{応用量子物性講義ノート}
\date{更新日\today}
\author{Yuto Masuda and Haruki Aoki}

\begin{document}
  \refe{hamiltonian-in-electro-magnetic-field-cm}より，電磁場中のハミルトニアンは，
  \begin{align}
    H = \frac{1}{2m}\qty(\bm{p} + e\bm{A})^2 - e\phi\label{hamiltonian-in-electro-magnetic-field-cm-re}
  \end{align}
  と書けるのであった．
  今回は$\phi = 0$とする．
  電磁場が十分弱いという条件のもと，\refe{hamiltonian-in-electro-magnetic-field-cm-re}を量子化すると，
  \begin{align}
    \hat{H} &= \frac{1}{2m}\qty(\hat{\bm{p}} + e\hat{\bm{A}})^2 \\
    &\simeq \frac{1}{2m}\qty(\hat{\bm{p}}^2 + e\hat{\bm{p}}\cdot\hat{\bm{A}} + e\hat{\bm{A}} \cdot \hat{\bm{p}})\\
  \end{align}
  と近似する．ここで，$\hat{H}^{(0)} \coloneqq \frac{\bm{p}^2}{2m}$，摂動項$\hat{V}(t) \coloneqq \frac{e}{2m} \qty(\hat{\bm{p}}\cdot\hat{\bm{A}} + \hat{\bm{A}} \cdot \hat{\bm{p}})$
  とする．さらに，$\div\bm{A} = 0$となるように$\bm{A}$を決める\footnote{Coulomb ゲージ}．すると，
  \begin{align}
    \qty(\hat{\bm{p}} \cdot \bm{A})\psi &= -\i\hbar \div\qty(\hat{\bm{A}} \psi)\\
    &= \i\hbar\qty[\qty(\div\hat{\bm{A}}) \psi + \hat{\bm{A}} \cdot \qty(\grad \psi)]\\
    &= \qty(\hat{\bm{A}} \cdot \hat{\bm{p}}) \psi
  \end{align}
  となる．よって，電子と電磁場の相互作用を摂動として加えたハミルトニアン，
  \begin{align}
    \hat{H} = \hat{H}^{(0)} + \frac{e}{m}\qty(\hat{\bm{A}} \cdot \hat{\bm{p}})
  \end{align}
  を得る．
  \begin{myex}{直線偏光}{}
    ベクトルポテンシャル$\hat{A}$が，
    \begin{align}
      \hat{\bm{A}}(\bm{r}, t) = 2A_0 \bm{e}_x \cos(\bm{k} \cdot \bm{r} - \omega t)
    \end{align}
    であるときを考える．
    ただし，$\bm{k} = \frac{\omega}{c}\bm{e}_z$とする．
    ベクトルポテンシャルと電磁場の関係より，
    \begin{align}
      \begin{cases}
      \bm{E}(\bm{r}, t) = \pdv{\hat{\bm{A}}}{t} = E_0 \bm{e}_x \sin(\bm{k} \cdot \bm{r} - \omega t)\\
      \bm{B}(\bm{r}, t) = \curl\bm{A} = \frac{E_0}{c}\bm{e}_y \sin(\bm{k} \cdot \bm{r} - \omega t)\\
      E_0 = -2 \omega A_0
      \end{cases}
    \end{align}
    が成り立っている．摂動項は，
    \begin{align}
      \hat{V}(t) &= \frac{e}{m}(\hat{\bm{A}}\cdot\hat{\bm{p}}) \\
      &= \frac{2eA_0}{m}\cos(\bm{k} \cdot \bm{r} - \omega t)\bm{e}_x\cdot\hat{\bm{p}} \\
      &= \frac{eA_0}{m}\qty[\e^{\i(\bm{k}\cdot\bm{r} - \omega t)} + \e^{-\i(\bm{k}\cdot\bm{r} - \omega t)}]\hat{p}_x\\
      &= \frac{eA_0}{m}\qty(\e^{\i\bm{k}\cdot\bm{r}}\hat{p}_x\e^{-\i\omega t} + \e^{-\i \bm{k}\cdot\bm{r}}\hat{p}_x \e^{\i \omega t})
    \end{align}
    と表せる．光の吸収を考えるときは，第1項$\frac{eA_0}{m}\e^{\i\bm{k}\cdot\bm{r}}\hat{p}_x$が支配的なのでこの項を$\hat{V}$とする．
    単位時間当たりの遷移確率を計算する．
    \refe{fermis-golden-rule}より，$\omega_{i\to f}$は，
    \begin{align}
      \omega_{i\to f} &= \frac{2\pi}{\hbar} \abs{\mel**{f}{\hat{V}}{i}}^2\delta(E_f - E_i - \hbar \omega) \\
      &= \frac{2 \pi}{\hbar}\qty(\frac{e A_0}{m})^2 \abs{\mel**{f}{\e^{\i\bm{k}\cdot \bm{r}}\hat{p}_x}{i}}^2 \delta(E_f - E_i - \hbar \omega)
    \end{align}
    と表せる．
    さて，$\abs{\mel**{f}{\e^{\i\bm{k}\cdot\bm{r}}\hat{p}_x}{i}}^2$を\textbf{電気双極子近似}を用いて計算する．
    電気双極子近似とは，電磁場の変化の項である$\e^{\i\bm{k}\cdot\bm{r}}$を$1$とみなす近似である．
    これは次のような議論から正当化される．
    原子の準位間隔は$E_f - E_i \sim 1\ \r{eV}$である．これと相互作用する電磁場のエネルギーは$\hbar \omega\sim 1\ \r{eV}$ である．これを波長に換算すると，
    \begin{align}
      \lambda = \frac{2\pi}{\hbar} = \frac{2\pi c}{\omega} \sim 1000\ \r{nm}
    \end{align}
    である．これは原子のスケール$1\ \si{\angstrom}$よりもはるかに大きいため，電子・原子を扱う上では電磁場は空間的に一様だとみなせる．よって，
    \begin{align}
      \e^{\i\bm{k}\cdot\bm{r}}\simeq 1 + \i \bm{k} \cdot \bm{r} + \cdots \simeq 1
    \end{align}
    と近似できる．
    \par
    電気双極子近似を用いると，
    \begin{align}
      \mel**{f}{\e^{\i \bm{k} \cdot \bm{r}}\hat{p}_x}{{i}} \simeq \mel**{f}{\hat{p}_x}{i}
    \end{align}
    を得る．さらに
    \begin{align}
      \qty[\hat{x}, \hat{H}^{(0)}] = \frac{\i \hbar}{m}\hat{p}_x
    \end{align}
    であるため
    \begin{align}
      \mel**{f}{\hat{p}_x}{i} &= \frac{m}{\i\hbar} \mel**{f}{\qty[\hat{x}, \hat{H}^{(0)}]}{i} \\ 
      &= \frac{m}{\i\hbar} \qty(\mel**{f}{\hat{x}\hat{H}^{(0)}}{i} - \mel**{f}{\hat{H}^{(0)}\hat{x}}{i}) \\ 
      &= \frac{m}{\i\hbar} \qty(\mel**{f}{\hat{x}E_i}{i} - \mel**{f}{E_f\hat{x}}{i}) \\ 
      &= \frac{m}{\i\hbar} (E_i - E_f)\mel**{f}{\hat{x}}{i}
    \end{align}
    を得る．よって，電磁場による単位時間当たりの遷移確率
    \begin{align}
      \omega_{i \to f} = \frac{2\pi}{\hbar^3} (eA_0)^2 (E_i - E_f)^2 \abs{\mel**{f}{\hat{x}}{i}}^2 \delta(E_f - E_i - \hbar\omega) 
    \end{align}
    を得る．これは$\mel**{f}{\hat{x}}{i} \neq 0$のときのみ$\omega_{i\to f} \neq 0$であるという選択則を表している．
  \end{myex}
\end{document}